\documentclass[11pt]{article}

\usepackage[T1]{fontenc}
\usepackage{lmodern}
\usepackage[margin=0.9in]{geometry}
\usepackage{microtype}
\usepackage{amsmath,amssymb,amsthm,mathtools}
\usepackage{booktabs}
\usepackage{enumitem}
\usepackage{xcolor}
\usepackage[colorlinks=true,linkcolor=blue!55!black,citecolor=blue!55!black,
  urlcolor=blue!55!black]{hyperref}

\newtheorem{theorem}{Theorem}[section]
\newtheorem{proposition}[theorem]{Proposition}
\newtheorem{lemma}[theorem]{Lemma}
\newtheorem{corollary}[theorem]{Corollary}
\theoremstyle{remark}
\newtheorem{remark}[theorem]{Remark}

\newcommand{\R}{\mathbb{R}}
\newcommand{\E}{\mathbb{E}}
\newcommand{\N}{\mathcal{N}}
\newcommand{\KL}{\operatorname{KL}}
\newcommand{\Ent}{\operatorname{Ent}}
\newcommand{\Law}{\operatorname{Law}}
\newcommand{\Id}{\operatorname{Id}}
\newcommand{\PiC}{\Pi}
\newcommand{\dd}{\,\mathrm{d}}
\newcommand{\push}{_{\#}}

\title{Entropic Transport Couplings for Gaussian Smoothing\\
and Variance-Preserving Diffusions}
\author{}
\date{}

\begin{document}
\maketitle

\begin{abstract}
This note gives a self-contained derivation of an exact entropic optimal
transport coupling associated with Gaussian smoothing and with the
variance-preserving Ornstein--Uhlenbeck diffusion used in diffusion models.
For the quadratic cost $c(x,y)=\frac12\lVert x-y\rVert^2$, an affine Gaussian
channel
\[
  Y=aX+\sqrt q\,Z, \qquad Z\sim\N(0,I_d),
\]
is the entropic optimizer between its input and output laws when the entropy
parameter is $\varepsilon=q/a$.  Gaussian smoothing corresponds to $a=1$,
while a variance-preserving diffusion has $q=1-a^2$.  We record the Sinkhorn
potentials, forward and reverse conditional laws, the score/barycentric
formula, the formulation with an OU-adapted cost, arbitrary two-time
couplings, and the infinite-time limit.
\end{abstract}

\tableofcontents

\section{Conventions and problem formulation}

Let $\mu,\nu\in\mathcal P_2(\R^d)$ be probability measures with finite
second moments.  Their set of couplings is denoted by
\[
  \PiC(\mu,\nu)
  :=\{\pi\in\mathcal P(\R^d\times\R^d):
       (P_1)\push\pi=\mu,\ (P_2)\push\pi=\nu\}.
\]
For $\varepsilon>0$, the entropic optimal transport (EOT) problem considered
throughout the note is
\begin{equation}
  \label{eq:eot}
  \inf_{\pi\in\PiC(\mu,\nu)}
  \left\{
    \int_{\R^d\times\R^d}\frac12\lVert x-y\rVert^2\,\dd\pi(x,y)
    +\varepsilon\KL(\pi\mid\mu\otimes\nu)
  \right\}.
\end{equation}
Here
\[
  \KL(\pi\mid r)
  =\begin{cases}
    \displaystyle\int\log\!\left(\frac{\dd\pi}{\dd r}\right)\dd\pi,
       &\pi\ll r,\\[1ex]
    +\infty,&\text{otherwise}.
  \end{cases}
\]
The normalization of the cost matters.  With the alternative cost
$\lVert x-y\rVert^2$, every matching value of $\varepsilon$ below must be
multiplied by $2$.

Under standard integrability assumptions, the minimizer of
\eqref{eq:eot} is unique and has the Sinkhorn form
\begin{equation}
  \label{eq:sinkhorn-form}
  \frac{\dd\pi^\star}{\dd(\mu\otimes\nu)}(x,y)
  =\exp\!\left(
    \frac{f(x)+g(y)-\frac12\lVert x-y\rVert^2}{\varepsilon}
  \right),
\end{equation}
where $f$ and $g$ are defined up to the gauge transformation
$(f,g)\mapsto(f+C,g-C)$.

\begin{remark}[Assumptions]
The formulas below are identities of measures and remain meaningful for a
singular input law $\mu$ (for instance, an empirical distribution).  To avoid
inessential qualifications when manipulating objective values, it is enough
to assume that the input has a finite second moment.  Convolution with a
nondegenerate Gaussian then gives a positive smooth output density and the
Gaussian channel has finite mutual information.  The endpoint $q=0$ is
understood as a weak limit.
\end{remark}

\section{The affine Gaussian channel theorem}

The following result contains both Gaussian smoothing and the VP diffusion
as special cases.

\begin{theorem}[Affine Gaussian channel as an EOT optimizer]
  \label{thm:affine}
Let $\mu\in\mathcal P_2(\R^d)$, let $a>0$ and $q>0$, and define
\begin{equation}
  \label{eq:affine-channel}
  X\sim\mu,\qquad Z\sim\N(0,I_d)\text{ independent},\qquad
  Y=aX+\sqrt q\,Z.
\end{equation}
Let $\nu=\Law(Y)$ and let $\rho$ be its Lebesgue density.  Then the joint law
\begin{equation}
  \label{eq:affine-coupling}
  \pi_{a,q}(\dd x,\dd y)
  =\mu(\dd x)\,\N(ax,qI_d)(\dd y)
\end{equation}
is the unique minimizer of \eqref{eq:eot} between $\mu$ and $\nu$ for
\begin{equation}
  \label{eq:epsilon-aq}
  \boxed{\ \varepsilon=\frac qa\ }.
\end{equation}
One choice of Sinkhorn potentials is
\begin{align}
  \label{eq:f-affine}
  f(x)&=\frac{1-a}{2}\lVert x\rVert^2,\\
  \label{eq:g-affine}
  g(y)&=-\frac{1-a}{2a}\lVert y\rVert^2
        -\varepsilon\log\rho(y)
        -\frac{d\varepsilon}{2}\log(2\pi q).
\end{align}
\end{theorem}

\begin{proof}
Because the second marginal of \eqref{eq:affine-coupling} has density $\rho$,
\begin{equation}
  \label{eq:rn-affine}
  \frac{\dd\pi_{a,q}}{\dd(\mu\otimes\nu)}(x,y)
  =\frac{(2\pi q)^{-d/2}}{\rho(y)}
    \exp\!\left(-\frac{\lVert y-ax\rVert^2}{2q}\right).
\end{equation}
The exponent is separated using
\begin{align}
  -\frac{\lVert y-ax\rVert^2}{2q}
  ={}&-\frac{\lVert x-y\rVert^2}{2\varepsilon}
     +\frac{a(1-a)}{2q}\lVert x\rVert^2
     -\frac{1-a}{2q}\lVert y\rVert^2,
  \label{eq:key-algebra}
\end{align}
where $\varepsilon=q/a$.  Substitution into \eqref{eq:rn-affine} gives
the Sinkhorn form \eqref{eq:sinkhorn-form} with
\eqref{eq:f-affine}--\eqref{eq:g-affine}.

For completeness, the optimality can also be seen directly.  If
$\pi\in\PiC(\mu,\nu)$, then, whenever the quantities are finite,
\begin{align}
 &\frac12\int\lVert x-y\rVert^2\dd\pi
  +\varepsilon\KL(\pi\mid\mu\otimes\nu) \notag\\
 &\qquad=\varepsilon\KL(\pi\mid\pi_{a,q})+C_{\mu,\nu,a,q},
  \label{eq:pythagorean}
\end{align}
where the constant is independent of $\pi$.  Indeed, insert
\eqref{eq:rn-affine}; all remaining terms depend only on the fixed first or
second marginal.  Relative entropy is nonnegative and vanishes only when
$\pi=\pi_{a,q}$.  A standard truncation argument extends the conclusion when
some intermediate terms are infinite.
\end{proof}

\begin{remark}[Why the matching is forced]
The cross term in the log-density of the channel is $a\langle x,y\rangle/q$.
The cross term in $-\lVert x-y\rVert^2/(2\varepsilon)$ is
$\langle x,y\rangle/\varepsilon$.  Hence the only possible matching for the
standard quadratic cost is $\varepsilon=q/a$.
\end{remark}

\subsection{An adapted-cost formulation}

Define
\begin{equation}
  \label{eq:adapted-cost-general}
  c_a(x,y):=\frac12\lVert y-ax\rVert^2.
\end{equation}
The same coupling is equivalently characterized by
\begin{equation}
  \label{eq:adapted-eot-general}
  \pi_{a,q}=\mathop{\rm argmin}_{\pi\in\PiC(\mu,\nu)}
  \left\{\int c_a\,\dd\pi
  q\KL(\pi\mid\mu\otimes\nu)\right\}.
\end{equation}
Indeed, its potentials for this problem can be chosen as
\begin{equation}
  \label{eq:adapted-potentials}
  \widetilde f(x)=0,
  \qquad
  \widetilde g(y)=-q\log\rho(y)-\frac{dq}{2}\log(2\pi q).
\end{equation}
This formulation makes the reference Gaussian channel completely explicit.

\section{Gaussian smoothing}

Let
\[
  \gamma_{\sigma^2}:=\N(0,\sigma^2I_d),
  \qquad \nu=\mu*\gamma_{\sigma^2}.
\]
Taking $a=1$ and $q=\sigma^2$ in Theorem~\ref{thm:affine} gives the following.

\begin{corollary}[Matched Gaussian smoothing]
  \label{cor:smoothing}
For the cost $\frac12\lVert x-y\rVert^2$, the EOT parameter matched to a
Gaussian smoothing of variance $\sigma^2$ is
\[
  \boxed{\ \varepsilon=\sigma^2\ }.
\]
The optimal coupling is the additive Gaussian coupling
\begin{equation}
  \label{eq:smoothing-coupling}
  \boxed{\quad
  \pi_\varepsilon(\dd x,\dd y)
  =\mu(\dd x)\,\gamma_\varepsilon(\dd y-x),
  \quad}
\end{equation}
or, probabilistically, $Y=X+\sqrt\varepsilon Z$.
\end{corollary}

If $\rho_\varepsilon$ denotes the density of
$\mu*\gamma_\varepsilon$, then
\begin{equation}
  \frac{\dd\pi_\varepsilon}
       {\dd(\mu\otimes(\mu*\gamma_\varepsilon))}(x,y)
  =\frac{(2\pi\varepsilon)^{-d/2}}{\rho_\varepsilon(y)}
   \exp\!\left(-\frac{\lVert x-y\rVert^2}{2\varepsilon}\right).
\end{equation}
Thus one may take
\begin{equation}
  f(x)=0,
  \qquad
  g(y)=-\varepsilon\log\rho_\varepsilon(y)
       -\frac{d\varepsilon}{2}\log(2\pi\varepsilon).
\end{equation}
The forward and reverse kernels are
\begin{align}
  Y\mid X=x&\sim\N(x,\varepsilon I_d),\\
  \pi_\varepsilon(\dd x\mid Y=y)
  &=\frac{\gamma_\varepsilon(y-x)\mu(\dd x)}
          {(\mu*\gamma_\varepsilon)(y)}.
\end{align}
Differentiating the Gaussian convolution yields Tweedie's identity
\begin{equation}
  \label{eq:tweedie-smoothing}
  \boxed{\quad
  \E[X\mid Y=y]
  =y+\varepsilon\nabla\log\rho_\varepsilon(y).
  \quad}
\end{equation}

If the heat semigroup is normalized as
$P_t\mu=\mu*\N(0,2tI_d)$, the matching is $\varepsilon=2t$.  If the cost is
$\lVert x-y\rVert^2$ instead of half that quantity, the matching becomes
$\varepsilon=2\sigma^2$.

\begin{remark}[Unmatched parameters]
For a general non-Gaussian $\mu$, the additive Gaussian coupling is not the
EOT optimizer when $\varepsilon\ne\sigma^2$ under our cost convention.  The
cross-term test in the preceding remark already shows the obstruction.  The
unmatched problem generally requires two nontrivial Sinkhorn potentials.
When both marginals are Gaussian, the optimizer is itself Gaussian and a
separate closed-form covariance calculation is possible.
\end{remark}

\section{Variance-preserving diffusion}

\subsection{The forward SDE and its marginals}

Let $\beta:[0,\infty)\to[0,\infty)$ be locally integrable and consider the
variance-preserving (VP) SDE
\begin{equation}
  \label{eq:vp-sde}
  \dd X_t=-\frac{\beta(t)}2X_t\,\dd t+\sqrt{\beta(t)}\,\dd W_t,
  \qquad X_0\sim\mu_0.
\end{equation}
Its Fokker--Planck equation is
\begin{equation}
  \label{eq:fokker-planck}
  \partial_t\rho_t
  =\frac{\beta(t)}2\nabla\!\cdot(x\rho_t)
   +\frac{\beta(t)}2\Delta\rho_t.
\end{equation}
After the time change $\tau(t)=\frac12\int_0^t\beta(r)\dd r$, this is the
Ornstein--Uhlenbeck equation associated with the quadratic potential
$V(x)=\frac12\lVert x\rVert^2$.

Define
\begin{equation}
  \label{eq:at-qt}
  B_t:=\int_0^t\beta(r)\dd r,
  \qquad
  a_t:=e^{-B_t/2},
  \qquad
  q_t:=1-a_t^2.
\end{equation}
Solving the linear SDE gives
\begin{equation}
  \label{eq:vp-solution}
  X_t=a_tX_0+\sqrt{q_t}\,Z,
  \qquad Z\sim\N(0,I_d)\text{ independent of }X_0,
\end{equation}
and hence
\begin{equation}
  \label{eq:vp-marginal}
  \mu_t=(a_t\Id)\push\mu_0*\N(0,q_tI_d).
\end{equation}
The identity $q_t=1-a_t^2$ is the variance-preserving property: if
$X_0\sim\N(0,I_d)$, then $X_t\sim\N(0,I_d)$ at every time.

\subsection{The exact finite-time Sinkhorn coupling}

Suppose $t>0$ and $B_t>0$, so $0<a_t<1$ and $q_t>0$.  The pathwise forward
coupling between $\mu_0$ and $\mu_t$ is
\begin{equation}
  \label{eq:vp-coupling}
  \boxed{\quad
  \pi_t(\dd x,\dd y)
  =\mu_0(\dd x)\,\N(a_tx,q_tI_d)(\dd y).
  \quad}
\end{equation}
Theorem~\ref{thm:affine}, with $a=a_t$ and $q=q_t$, shows that this is exactly
the standard quadratic EOT coupling for
\begin{equation}
  \label{eq:epsilon-t}
  \boxed{\quad
  \varepsilon_t=\frac{q_t}{a_t}
  =\frac{1-a_t^2}{a_t}
  =a_t^{-1}-a_t
  =2\sinh\!\left(\frac{B_t}{2}\right).
  \quad}
\end{equation}
Writing $\rho_t$ for the density of $\mu_t$, its density ratio is
\begin{equation}
  \label{eq:vp-rn}
  \frac{\dd\pi_t}{\dd(\mu_0\otimes\mu_t)}(x,y)
  =\frac{(2\pi q_t)^{-d/2}}{\rho_t(y)}
    \exp\!\left(-\frac{\lVert y-a_tx\rVert^2}{2q_t}\right),
\end{equation}
and one may choose the potentials
\begin{align}
  \label{eq:vp-f}
  f_t(x)&=\frac{1-a_t}{2}\lVert x\rVert^2,\\
  \label{eq:vp-g}
  g_t(y)&=-\frac{1-a_t}{2a_t}\lVert y\rVert^2
          -\varepsilon_t\log\rho_t(y)
          -\frac{d\varepsilon_t}{2}\log(2\pi q_t).
\end{align}

At $t=0$, $a_0=1$, $q_0=0$, and $\varepsilon_0=0$.  The coupling is the
diagonal law $(\Id,\Id)\push\mu_0$.  This endpoint is a zero-noise limit, not
an entropy-regularized problem with positive $\varepsilon$.

\subsection{OU-adapted cost}

For the time-dependent cost
\begin{equation}
  \label{eq:ou-cost}
  c_t^{\mathrm{OU}}(x,y)
  :=\frac12\lVert y-a_tx\rVert^2,
\end{equation}
the same coupling satisfies
\begin{equation}
  \label{eq:ou-eot}
  \boxed{\quad
  \pi_t=\mathop{\rm argmin}_{\pi\in\PiC(\mu_0,\mu_t)}
  \left\{
    \int c_t^{\mathrm{OU}}\dd\pi
    q_t\KL(\pi\mid\mu_0\otimes\mu_t)
  \right\}.
  \quad}
\end{equation}
The potentials simplify to
\begin{equation}
  \widetilde f_t(x)=0,
  \qquad
  \widetilde g_t(y)
  =-q_t\log\rho_t(y)-\frac{dq_t}{2}\log(2\pi q_t).
\end{equation}
This is the most direct one-step Schr\"odinger formulation: the Gibbs kernel
is precisely the OU transition kernel.  Since that reference process is
initialized at $\mu_0$, its endpoint law is already $\mu_t$, and its joint
endpoint law is therefore feasible and minimizes the relative entropy.

\subsection{Reverse kernel and denoising formula}

The forward conditional law is
\[
  X_t\mid X_0=x\sim\N(a_tx,q_tI_d).
\]
Bayes' formula gives the reverse conditional law
\begin{equation}
  \label{eq:reverse-kernel}
  \boxed{\quad
  \pi_t(\dd x\mid X_t=y)
  =\frac{
    \exp\!\left(-\frac{\lVert y-a_tx\rVert^2}{2q_t}\right)\mu_0(\dd x)
  }{
    \displaystyle\int
    \exp\!\left(-\frac{\lVert y-a_tz\rVert^2}{2q_t}\right)\mu_0(\dd z)
  }.
  \quad}
\end{equation}
Differentiation under the integral in \eqref{eq:vp-marginal} yields
\begin{equation}
  \nabla\log\rho_t(y)
  =-\frac{1}{q_t}\bigl(y-a_t\E[X_0\mid X_t=y]\bigr).
\end{equation}
Consequently, the reverse barycentric projection is
\begin{equation}
  \label{eq:vp-tweedie}
  \boxed{\quad
  \E[X_0\mid X_t=y]
  =\frac1{a_t}\left(y+q_t\nabla\log\rho_t(y)\right).
  \quad}
\end{equation}
This is the VP version of Tweedie's formula and is the precise relation
between the entropic coupling, denoising, and the score of the noised data
distribution.

\subsection{Any two diffusion times}

The same result holds between arbitrary times $0\le s<t$.  Set
\begin{equation}
  \label{eq:ast-qst}
  a_{s,t}:=\exp\!\left(-\frac12\int_s^t\beta(r)\dd r\right),
  \qquad
  q_{s,t}:=1-a_{s,t}^2.
\end{equation}
The Markov transition is
\[
  X_t=a_{s,t}X_s+\sqrt{q_{s,t}}\,Z_{s,t},
  \qquad Z_{s,t}\sim\N(0,I_d),\quad Z_{s,t}\perp X_s.
\]
Therefore
\begin{equation}
  \label{eq:two-time-coupling}
  \pi_{s,t}(\dd x,\dd y)
  =\mu_s(\dd x)\,\N(a_{s,t}x,q_{s,t}I_d)(\dd y)
\end{equation}
is the exact standard quadratic Sinkhorn coupling between $\mu_s$ and $\mu_t$
with
\begin{equation}
  \label{eq:epsilon-st}
  \boxed{\quad
  \varepsilon_{s,t}=\frac{q_{s,t}}{a_{s,t}}
  =2\sinh\!\left(\frac12\int_s^t\beta(r)\dd r\right).
  \quad}
\end{equation}
All preceding formulas apply after replacing
$(\mu_0,\mu_t,a_t,q_t)$ with
$(\mu_s,\mu_t,a_{s,t},q_{s,t})$.

\subsection{Common time normalizations}

If $\beta(t)\equiv\beta$ is constant, then
\[
  a_t=e^{-\beta t/2},
  \qquad q_t=1-e^{-\beta t},
  \qquad \varepsilon_t=2\sinh(\beta t/2).
\]
For the common OU normalization
\begin{equation}
  \dd X_t=-X_t\dd t+\sqrt2\,\dd W_t,
\end{equation}
one has $\beta=2$ and hence
\begin{equation}
  a_t=e^{-t},
  \qquad q_t=1-e^{-2t},
  \qquad \boxed{\varepsilon_t=2\sinh t}.
\end{equation}

\section{The infinite-time limit}

Assume
\begin{equation}
  \label{eq:infinite-noise}
  \int_0^\infty\beta(r)\dd r=+\infty.
\end{equation}
Then $a_t\to0$, $q_t\to1$, and
\begin{equation}
  \mu_t\Longrightarrow\gamma:=\N(0,I_d).
\end{equation}
The forward endpoint coupling converges weakly to
\begin{equation}
  \label{eq:infinite-coupling}
  \boxed{\quad
  \pi_\infty(\dd x,\dd y)
  =\mu_0(\dd x)\gamma(\dd y)
  =\mu_0\otimes\gamma.
  \quad}
\end{equation}
Thus the initial state and the stationary endpoint are independent.

For the standard quadratic cost,
\begin{equation}
  \varepsilon_t=\frac{1-a_t^2}{a_t}\longrightarrow+\infty.
\end{equation}
Accordingly, \eqref{eq:infinite-coupling} is the infinite-regularization limit
of the standard quadratic Sinkhorn coupling.  Except in degenerate cases
(for example, a deterministic first marginal), the product coupling is not
the optimizer for any finite $\varepsilon$: at finite regularization the
quadratic Gibbs kernel retains a nonzero $x$--$y$ interaction.

For the OU-adapted formulation,
\[
  c_\infty^{\mathrm{OU}}(x,y)=\frac12\lVert y\rVert^2,
  \qquad q_\infty=1.
\]
The cost depends only on the fixed second marginal and is therefore constant
over $\PiC(\mu_0,\gamma)$.  It follows directly that
\begin{equation}
  \pi_\infty
  =\mathop{\rm argmin}_{\pi\in\PiC(\mu_0,\gamma)}
    \KL(\pi\mid\mu_0\otimes\gamma)
  =\mu_0\otimes\gamma.
\end{equation}
Unlike the standard-cost parameter, the adapted entropy parameter has the
finite limit $q_\infty=1$; the interaction instead disappears from the cost
because $a_t\to0$.

\section{Objective value and mutual information}

The KL term of the forward coupling has a direct information-theoretic
meaning:
\begin{equation}
  \KL(\pi_{a,q}\mid\mu\otimes\nu)=I(X;Y).
\end{equation}
Moreover,
\begin{equation}
  \E\lVert X-Y\rVert^2
  =(1-a)^2\E\lVert X\rVert^2+dq.
\end{equation}
Hence the minimum value of the standard-cost problem in
Theorem~\ref{thm:affine} is
\begin{equation}
  \label{eq:optimal-value}
  \frac12\left((1-a)^2\E\lVert X\rVert^2+dq\right)
  +\frac qa I(X;Y).
\end{equation}
When differential entropies are finite,
\begin{equation}
  I(X;Y)=h(Y)-h(Y\mid X)
  =h(\nu)-\frac d2\log(2\pi e q).
\end{equation}
For the VP diffusion, substitute $(a,q)=(a_t,q_t)$ in
\eqref{eq:optimal-value}.  As $t\to\infty$, the mutual information tends to
zero under standard moment assumptions, consistently with the limiting
product coupling.

\section{Summary of parameter conventions}

\begin{table}[ht]
  \centering
  \small
  \begin{tabular}{@{}lll@{}}
    \toprule
    Setting & Cost and entropy term & Matching parameter \\
    \midrule
    Gaussian smoothing
      & $\frac12\lVert x-y\rVert^2+\varepsilon\KL$
      & $\varepsilon=\sigma^2$ \\
    Gaussian smoothing
      & $\lVert x-y\rVert^2+\varepsilon\KL$
      & $\varepsilon=2\sigma^2$ \\
    VP, finite $t$
      & $\frac12\lVert x-y\rVert^2+\varepsilon_t\KL$
      & $\varepsilon_t=(1-a_t^2)/a_t$ \\
    VP, finite $t$
      & $\frac12\lVert y-a_tx\rVert^2+q_t\KL$
      & $q_t=1-a_t^2$ \\
    VP, $t=\infty$
      & standard quadratic cost
      & $\varepsilon_t\to+\infty$ \\
    VP, $t=\infty$
      & $\frac12\lVert y\rVert^2+\KL$
      & product coupling \\
    \bottomrule
  \end{tabular}
  \caption{The entropy coefficient under the cost convention used in this
  note.  Here $\KL$ abbreviates $\KL(\pi\mid\text{product of marginals})$.}
  \label{tab:summary}
\end{table}

The central identities are therefore
\begin{equation}
  \boxed{
  \begin{gathered}
  X_t=a_tX_0+\sqrt{1-a_t^2}\,Z,\\
  \pi_t=\Law(X_0,X_t),\\
  \varepsilon_t=\frac{1-a_t^2}{a_t},\\
  \E[X_0\mid X_t=y]
  =a_t^{-1}\bigl(y+(1-a_t^2)\nabla\log\rho_t(y)\bigr),\\
  \pi_t\Longrightarrow\mu_0\otimes\N(0,I_d)
  \quad\text{as }t\to\infty.
  \end{gathered}}
\end{equation}

\begin{thebibliography}{9}
\small
\setlength{\itemsep}{0.35\baselineskip}

\bibitem{peyre-cuturi}
G. Peyr\'e and M. Cuturi,
\emph{Computational Optimal Transport},
Foundations and Trends in Machine Learning, 2019.

\bibitem{leonard}
C. L\'eonard,
\emph{A survey of the Schr\"odinger problem and some of its connections with
optimal transport},
Discrete and Continuous Dynamical Systems A, 2014.

\bibitem{ho}
J. Ho, A. Jain, and P. Abbeel,
\emph{Denoising Diffusion Probabilistic Models},
Advances in Neural Information Processing Systems, 2020.

\bibitem{song}
Y. Song, J. Sohl-Dickstein, D. P. Kingma, A. Kumar, S. Ermon, and B. Poole,
\emph{Score-Based Generative Modeling through Stochastic Differential
Equations},
International Conference on Learning Representations, 2021.

\end{thebibliography}

\end{document}
